\section*{Chapitre \ref{ch:rs:order-book-features} --- Variables du carnet d'ordres}

\begin{solution}{exo:rs:order-book-features:1}
$I = (700 - 300)/1\,000 = 0.4$~; le prix milieu vaut 99.985~; le \omterm{def:rs:order-book-features:micro}{mid pondéré} vaut $(99.99 \times 700 + 99.98 \times 300)/1\,000 = 99.987$, déplacé vers
le cours vendeur, parce que sa file est la plus petite.
\end{solution}

\begin{solution}{exo:rs:order-book-features:2}
Le cours acheteur a augmenté, donc $+q^b_n = +200$ et l'ancienne file acheteuse n'est pas retranchée~; le cours vendeur n'a pas changé, donc $-400 + 400 = 0$~: $e_n = +200$.
\end{solution}

\begin{solution}{exo:rs:order-book-features:3}
$11\,000/2\,750 = 4$ tailles moyennes de file~; $4 \times 0.47 = 1.88$ pas de cotation.
\end{solution}

\begin{solution}{exo:rs:order-book-features:4}
Avec $\lambda = \mu$, une file de $n$ lots met un temps à se vider qui croît sans borne en probabilité lorsque $n$ croît, tandis que la file à un seul lot
se vide en un temps qui ne dépend pas de $n$~; ainsi le cours vendeur survit presque sûrement au cours acheteur, et P(hausse) tend vers zéro. Les classes extrêmes du modèle
sont donc proches de 0 et de 1, alors que les données restent entre 0.25 et 0.74~: quelque chose d'autre que les marches aléatoires propres aux files fait
bouger le prix.
\end{solution}

\begin{solution}{exo:rs:order-book-features:5}
À une seconde, la plupart des intervalles n'ont aucune variation du prix milieu, et une variation est un seul tick déclenché par un événement dont l'OFI ne voit qu'une partie~; sur
des intervalles plus longs, les variations sont des sommes de nombreux événements, leur caractère discret s'estompe en moyenne, et la relation linéaire avec le flux cumulé apparaît.
\end{solution}

\begin{solution}{exo:rs:order-book-features:6}
$(a q^b + b q^a)/(q^b + q^a) = \frac{a + b}{2} + \frac{a - b}{2}\cdot\frac{q^b - q^a}{q^b + q^a} = m + \frac s2 I$.
\end{solution}

\begin{solution}{exo:rs:order-book-features:7}
0.277 pour le \omterm{def:rs:order-book-features:depth}{déséquilibre de profondeur} à cinq niveaux contre 0.253 pour le déséquilibre au meilleur niveau~: les niveaux plus profonds apportent ici de l'information parce que
les fournisseurs de liquidité du simulateur annulent aussi les cotations périmées derrière le meilleur niveau. Sur des données réelles, les niveaux profonds sont aussi ceux où les ordres sont
les moins coûteux à falsifier.
\end{solution}

\begin{solution}{exo:rs:order-book-features:8}
Les variables doivent être calculées à l'instant où la société pouvait les obtenir, c'est-à-dire l'heure de réception, et non l'horodatage de la bourse~; un carnet consolidé construit
à partir des horodatages de la bourse suppose que les messages de toutes les plateformes arrivent en même temps, ce qui n'est pas le cas. La précision est aussi peu informative sans le
taux de base~: aux déséquilibres extrêmes, la direction naïve a déjà raison trois fois sur quatre. Recalculer avec des carnets à l'heure de réception et présenter la
précision par classe, par rapport au taux de base.
\end{solution}

\begin{solution}{pb:rs:order-book-features:1}
\textbf{1.} 0.245 et 0.736. \textbf{2.} 16\,564 et 14\,538. \textbf{3.} 3.4 contre 55.4 lots, et 49.2 contre 3.1. \textbf{4.} Proche de zéro,
légèrement en dessous (la classe centrée sur $-0.1$ donne 0.500). \textbf{5.} 2.03 lots par seconde arrivent et 2.13 sortent. \textbf{6.} 0.017 et 0.979.
\textbf{7.} Il traite les files comme des marches aléatoires indépendantes, alors que le prix efficient détermine quel côté est annulé et contre lequel on négocie,
et une nouvelle liquidité continue d'arriver à la plus petite file. \textbf{8.} Le prix efficient~: les ordres informés et les annulations de cotations périmées qui
dépendent de sa position. \textbf{9.} 10~\%, 53~\% et 79~\%~; 0.47 pas de cotation par taille moyenne de la meilleure file. \textbf{10.} 2\,747 actions (27.5 lots).
\textbf{11.} Une erreur quadratique moyenne inférieure de 10~\% à celle du prix milieu. \textbf{12.} Le \omterm{def:rs:order-book-features:micro}{mid pondéré}~: 12~\% et 3~\% en dessous du prix milieu. \textbf{13.} 0.30 et
0.25. \textbf{14.} Parce qu'un modèle étudié sur un jeu de chiffres et négocié sur un autre est un autre modèle~; tout écart est un
bogue silencieux. \textbf{15.} 2\,766 lignes bilatérales sur 2\,769 messages. \textbf{16.} L'OFI et les taux ont besoin de chaque événement dans l'ordre~; les déséquilibres,
le \omterm{def:rs:order-book-features:micro}{mid pondéré} et le \omterm{def:rs:order-book-features:micro}{microprix} peuvent être calculés à partir d'instantanés du carnet. \textbf{17.} La profondeur au-delà du meilleur niveau, et le déséquilibre au meilleur niveau lui-même~:
les ordres passés sans intention de négocier. \textbf{18.} \emph{Résultat nommé~:} par dixième de déséquilibre de file, la probabilité d'un mouvement à la hausse varie
de 0.245 à 0.736 sur la bande simulée, là où la course de deux files de naissance et de mort prédit 0.017 à 0.979~; l'OFI explique 53~\% des variations à dix
secondes du prix milieu. \textbf{19.} Le prix efficient~: des arrivées et des annulations qui en dépendent (un modèle à files réactives avec un état caché).
\textbf{20.} L'état de l'offre et de la demande à l'instant présent et, avec son flux, là où le prix est sur le point d'aller.
\end{solution}

\begin{solution}{iq:rs:order-book-features:1}
Plutôt à la hausse~: la petite file du cours vendeur est bien plus proche de se vider, et les fournisseurs de liquidité côté cours vendeur sont peut-être ceux qui se retirent. Mais pas avec
la certitude qu'une course naïve entre files laisse supposer~; sur la bande simulée, le dixième le plus chargé côté acheteur donne environ 0.74, et non 0.98.
\iqlookfor{direction et calibration~: le déséquilibre est informatif mais bruité.}
\end{solution}

\begin{solution}{iq:rs:order-book-features:2}
Quantité nette ajoutée au cours acheteur, moins celle retirée de ce cours, moins celle ajoutée au cours vendeur, plus celle retirée de ce dernier, en comptant des files entières lorsque les prix changent.
Un flux d'ordres net de taille $x$ consomme ou construit des files de taille typique $D$~; le prix se déplace d'environ un pas de cotation par file consommée, de sorte que la
variation est proportionnelle à $x/D$~: linéaire, avec une pente inversement proportionnelle à la profondeur.
\iqlookfor{les quatre termes et l'argument de profondeur.}
\end{solution}

\begin{solution}{iq:rs:order-book-features:3}
Une estimation du prix équitable meilleure que le prix milieu, à partir du déséquilibre (et de la fourchette)~: l'espérance du prix milieu futur sachant le carnet d'ordres. Estimer par
découpage du déséquilibre (et de la fourchette) en classes et moyenne des variations ultérieures du prix milieu, en un pas ou itérée jusqu'à convergence, sur des données passées~; valider hors
échantillon, par rapport au prix milieu et au \omterm{def:rs:order-book-features:micro}{mid pondéré}.
\iqlookfor{conditionnement à l'état, estimation hors échantillon, la comparaison.}
\end{solution}

\begin{solution}{iq:rs:order-book-features:4}
Une implémentation de référence et un jeu d'essai commun~: rejouer les mêmes messages enregistrés dans chaque implémentation et exiger des
sorties identiques (au dernier bit lorsque c'est possible, à une tolérance serrée sinon) en intégration continue, avec le même ordre d'opérations dans
l'arithmétique.
\iqlookfor{un test de rejeu sur un jeu d'essai commun, et non une revue de code.}
\end{solution}

\begin{solution}{iq:rs:order-book-features:5}
La taille affichée puis retirée sans intention de négocier fausse les variables de déséquilibre et de profondeur (spoofing, layering)~; les variables de flux comptent ces ajouts
comme de la demande. Robustesse~: pondérer les quantités selon leur durée de présence ou la probabilité qu'elles soient exécutées, réduire le poids des niveaux au-delà du meilleur niveau, utiliser
le flux exécuté en plus de la taille cotée, et surveiller les suites d'ajouts et d'annulations de grande taille.
\iqlookfor{le mécanisme et au moins deux parades.}
\end{solution}

\begin{solution}{iq:rs:order-book-features:6}
P(hausse) est la probabilité que la file du cours vendeur se vide en premier, croissante en la taille côté acheteur et décroissante en la taille côté vendeur~; à taux égaux, elle vaut
un demi à tailles égales. Elle laisse de côté le prix efficient (flux informé et annulations de cotations périmées qui rendent les épuisements dépendants de
lui), les taux réactifs aux files, et les nouveaux niveaux à l'intérieur de la fourchette~; elle est donc trop confiante aux déséquilibres extrêmes.
\iqlookfor{la course, sa monotonie, et ce qui la rend trop confiante.}
\end{solution}
